\chapter{Volatility Targeting and Risk-Managed Portfolios}\label{ch:s1:volatility-targeting-and-risk-managed-portfolios}

When volatility jumps, every portfolio that sizes itself to a volatility target must sell, the same day, for the same reason. Moreira and Muir found
that taking less risk when volatility is high raised the Sharpe ratios of the market and of most equity factors; Cederburg, O'Doherty, Wang and Yan,
testing 103 strategies as a real-time investor would, found that \omterm{def:s1:volatility-targeting-and-risk-managed-portfolios:managed}{volatility-managed portfolios} do not systematically beat the unmanaged ones. On this
chapter's synthetic markets both are true: managing the market by its volatility improves the Sharpe ratio by 0.03 on average over twenty
ten-year samples, and does better in only half of them. What is certain is the selling: after a synthetic 6\% fall, funds holding a fifth of a percent
of the market sell 14.5\% of the next day's volume. The build is \texttt{firm.voltarget}.

\section{The mechanics of volatility targeting}

Chapter~19 defined \omterm{def:s1:trend-following:voltarget}{volatility targeting}: a position sized at a target volatility divided by its forecast. A risk-managed portfolio adds a rule for the
forecast and for how often to trade.

\begin{definition}[Volatility-managed portfolio]\label{def:s1:volatility-targeting-and-risk-managed-portfolios:managed}
A \emph{volatility-managed portfolio}\index{volatility-managed portfolio} scales an underlying portfolio's exposure by the inverse of its forecast
volatility (or variance), so that it takes less risk after volatile periods and more after calm ones, usually with a cap on leverage and a band within
which the exposure is not rebalanced.
\end{definition}

The ingredients are few (\cref{lst:s1:volatility-targeting-and-risk-managed-portfolios:weights}): a variance forecast (the previous month's realised
variance, or an exponentially weighted one as in Book~4, chapter~18); a scaling rule, target over volatility, or Moreira and Muir's target squared over
variance, which moves exposure faster; a leverage cap; a rebalancing band. Each forecast rule and band gives a different turnover, and the turnover is
the strategy's cost.

\section{Does it improve the Sharpe ratio?}

Moreira and Muir's argument is simple: if expected returns do not rise in proportion to volatility, scaling exposure down when volatility is high
removes risk faster than it removes return. \texttt{firm.synthmkt}'s market factor satisfies that by construction: a constant 6\% premium, GJR-GARCH
volatility with a leverage effect. Twenty ten-year samples of it, held at a constant exposure or managed by the previous month's realised variance
(target 16\%, leverage capped at three):

\begin{center}
\small
\begin{tabular}{@{}lccc@{}}
\toprule
twenty ten-year samples & constant exposure & inverse volatility & inverse variance\\
\midrule
mean Sharpe ratio & 0.49 & 0.52 & 0.48\\
samples better than constant & --- & 10 & 8\\
mean largest drawdown (annual volatilities) & 2.70 & 2.58 & 2.67\\
\bottomrule
\end{tabular}
\end{center}

The mean gain of the inverse-volatility portfolio is 0.028, with a standard deviation of 0.22 across samples
(\cref{fig:s1:volatility-targeting-and-risk-managed-portfolios:seeds}). The true Sharpe ratio of the constant portfolio is $0.06/0.16 = 0.375$; the
samples range from $-0.44$ to 1.18. A gain of a few hundredths is real in the model and invisible in any one decade, and the more aggressive inverse
variance loses it to the noise of its own forecasts. That is Cederburg and co-authors' finding in a laboratory: the improvement exists in the
population, and a real-time investor with one history cannot count on it.

\begin{omfigure}
\begin{tikzpicture}
\begin{axis}[width=8cm,height=7cm,xlabel={\small Sharpe ratio, constant exposure},ylabel={\small Sharpe ratio, volatility-managed},
  tick label style={font=\footnotesize},xmin=-0.6,xmax=1.4,ymin=-0.6,ymax=1.4,grid=major,grid style={gray!25},table/col sep=comma]
\addplot[only marks,mark=*,mark size=1.8pt,omThm] table[x=constant,y=managed]{figdata/strategies-1/26-volatility-targeting-and-risk-managed-portfolios/seeds.csv};
\addplot[black!50,dashed] coordinates {(-0.6,-0.6) (1.4,1.4)};
\end{axis}
\end{tikzpicture}

\omcaption{Twenty ten-year samples of the synthetic market factor (GJR-GARCH, constant premium): the Sharpe ratio of the volatility-managed market
(inverse of the previous month's realised volatility, target 16\%, leverage capped at three) against that of the constant exposure; points above the
dashed line are samples where management helped. Data: \texttt{s1\_voltarget.seeds}.}\label{fig:s1:volatility-targeting-and-risk-managed-portfolios:seeds}
\end{omfigure}

Harvey, Hoyle, Korgaonkar, Rattray, Sargaison and van Hemert found where the gain lives: only in risk assets such as equities and credit, where
volatility rises when prices fall (the leverage effect); for bonds, currencies and commodities the Sharpe ratio hardly changes. What targeting does
everywhere is shrink the tails, because the worst days tend to come when volatility is already high and the exposure already small.
\texttt{firm.synthfut}'s classes have volatility clustering but no leverage effect, and targeting them at 10\% changes their Sharpe ratios in both
directions: 0.19 to 0.19 for equities, 0.35 to 0.28 for bonds, $-0.18$ to $-0.08$ for currencies, $-0.21$ to $-0.25$ for commodities. Their largest
drawdowns fall for bonds (6.19 to 4.99 annual volatilities) and currencies (13.28 to 11.45) and rise for equities and commodities.

\section{The industry's footprint}

\begin{definition}[Leverage rebalancing flow]\label{def:s1:volatility-targeting-and-risk-managed-portfolios:flow}
A \emph{leverage rebalancing flow}\index{leverage rebalancing flow} is the net trading of portfolios that follow mechanical exposure rules
(volatility targets, leveraged funds' daily resets, risk parity) in response to price and volatility changes; because many follow similar rules, the
flow is concentrated and predictable from the rules and the market's recent moves.
\end{definition}

\texttt{firm.voltarget}'s flow simulator (\cref{lst:s1:volatility-targeting-and-risk-managed-portfolios:flow}) puts funds holding 0.2\% of a market's
value into a market that trades 0.4\% of its value a day. They target 10\% volatility with an exponentially weighted forecast (centre of mass 20 days,
leverage capped at two). The market is calm at 12\% for 250 days, falls 6\% in a day, and runs at 40\% for twenty days. The funds start at an exposure
of 0.73; the day after the fall their forecast jumps and they sell 14.5\% of the day's volume, and over the twenty days a total of 20.5\% of one day's
volume, ending at an exposure of 0.25 (\cref{fig:s1:volatility-targeting-and-risk-managed-portfolios:spike}). All of the flow parameters are assumptions:
the size of volatility-targeting money is not measured by any source fetched for this chapter.

\begin{omfigure}
\begin{tikzpicture}
\begin{axis}[width=11cm,height=5cm,xlabel={\small days from the 6\% fall},ylabel={\small flow (\% of the day's volume)},
  tick label style={font=\footnotesize},xmin=-10,xmax=49,ymin=-16,ymax=2,grid=major,grid style={gray!25},axis y line*=left,
  table/col sep=comma]
\addplot[ybar,bar width=3pt,fill=omThm!60,draw=omThm] table[x=day,y=flow_pct]{figdata/strategies-1/26-volatility-targeting-and-risk-managed-portfolios/spike.csv};
\end{axis}
\begin{axis}[width=11cm,height=5cm,xmin=-10,xmax=49,ymin=0,ymax=1,axis y line*=right,axis x line=none,ylabel={\small exposure},
  tick label style={font=\footnotesize},table/col sep=comma]
\addplot[omDef,thick] table[x=day,y=exposure]{figdata/strategies-1/26-volatility-targeting-and-risk-managed-portfolios/spike.csv};
\end{axis}
\end{tikzpicture}

\omcaption{Volatility-targeting funds around a synthetic 6\% fall: their net flow as a share of each day's volume (bars, left) and their exposure
(line, right). Funds hold 0.2\% of the market's value and the market trades 0.4\% of its value a day (assumed). Data:
\texttt{s1\_voltarget.spike}.}\label{fig:s1:volatility-targeting-and-risk-managed-portfolios:spike}
\end{omfigure}

The flow moves prices, and the moves feed the forecasts. With a linear impact of 0.1 times the flow's share of volume, the twenty days' fall deepens
from 16.6\% to 18.7\%; with 0.3, to 22.0\%. The selling is predictable from public information (the market's own volatility), which is what makes it
a strategy file for someone else: whoever provides liquidity to deleveraging funds is paid for it.

\section{Strategy files}

\begin{strategyfile}{Volatility-managed equity}\label{strat:s1:volatility-targeting-and-risk-managed-portfolios:equity}
\sfield{Who pays you, and why} Nobody in particular: expected returns do not rise as much as volatility, so scaling down in volatile periods improves the
risk-return trade-off, if it does.
\sfield{Instruments and venues} Equity index futures; equity factor portfolios.
\sfield{Signal} Forecast volatility (realised or EWMA).
\sfield{Sizing and execution} Exposure at target over forecast, capped; bands to limit turnover.
\sfield{Costs} Turnover from forecast changes, highest with inverse variance.
\sfield{How it dies} Forecast noise; instability of the relation between volatility and returns.
\sfield{Horizon, capacity, infrastructure} Monthly or daily rebalancing; large capacity.
\sfield{Backtest honestly} Real-time forecasts and scaling constants; compare directly with the unmanaged portfolio, not only by alpha.
\sfield{Sources} Moreira and Muir (2017); Cederburg, O'Doherty, Wang and Yan (2020); this chapter: $+0.03$ on average, better in 10 of 20 samples.
\end{strategyfile}

\begin{strategyfile}{Volatility-targeted multi-asset}\label{strat:s1:volatility-targeting-and-risk-managed-portfolios:multi}
\sfield{Who pays you, and why} The asset classes' premia, at a steadier risk.
\sfield{Instruments and venues} Futures in equities, bonds, credit, commodities.
\sfield{Signal} Each asset's and the portfolio's forecast volatility.
\sfield{Sizing and execution} Scaling at asset and portfolio level; leverage cap.
\sfield{Costs} Moderate.
\sfield{How it dies} Correlations that jump with volatility, so that diversification disappears when the target cuts exposure.
\sfield{Horizon, capacity, infrastructure} Daily to monthly; a risk system.
\sfield{Backtest honestly} Forecasts from past data only; tail statistics as well as Sharpe ratios.
\sfield{Sources} Harvey, Hoyle, Korgaonkar, Rattray, Sargaison and van Hemert (2018).
\end{strategyfile}

\begin{strategyfile}{Trading against deleveraging flow}\label{strat:s1:volatility-targeting-and-risk-managed-portfolios:against}
\sfield{Who pays you, and why} Volatility-targeting and other mechanical funds that must sell after volatility jumps.
\sfield{Instruments and venues} Index futures near the close, when the funds rebalance.
\sfield{Signal} The forecast selling: the funds' rules applied to the market's recent returns and an estimate of their size.
\sfield{Sizing and execution} Buy into the forecast selling, exit as the pressure fades; small against the crash risk.
\sfield{Costs} Low; the risk is the continuation of the fall.
\sfield{How it dies} Falls that continue for fundamental reasons; competition from others who read the same flows.
\sfield{Horizon, capacity, infrastructure} Days; an estimate of the funds' assets and rules.
\sfield{Backtest honestly} Only information available at the time; the crash's continuation in the sample.
\sfield{Sources} No performance figure verified; this chapter's flow model.
\end{strategyfile}

\section{Tutorial: mechanical selling}\label{tut:s1:volatility-targeting-and-risk-managed-portfolios:tut}

\begin{tutorial}
\textbf{Goal.} Manage the synthetic market's volatility over twenty samples and four asset classes, and simulate the selling of targeting funds after a
spike, with and without feedback. \textbf{End state:} the table and the two figures.
\begin{tutsteps}
\item \textbf{Forecasts and weights}.
\omcode{code/firm/voltarget/firm_voltarget.py}{26}{48}{Variance forecasts and target weights.}\label{lst:s1:volatility-targeting-and-risk-managed-portfolios:weights}
\item \textbf{The aggregate flow}: funds targeting volatility, their trade as a share of volume, and its impact fed back.
\omcode{code/firm/voltarget/firm_voltarget.py}{70}{91}{Targeting funds' flow and its feedback.}\label{lst:s1:volatility-targeting-and-risk-managed-portfolios:flow}
\item \textbf{Run} \texttt{summary()}, \texttt{classes()} and \texttt{spike(impact)} for three impacts, and \texttt{fig\_voltarget.py}.
\end{tutsteps}
\textbf{What to change next.} Add a rebalancing band and measure its effect on turnover and on the Sharpe gain; let several fund types (volatility
targets, leveraged ETFs' daily resets) trade together; let the size of the funds grow over time.
\end{tutorial}

\section{Build: volatility targeting}\label{bld:s1:volatility-targeting-and-risk-managed-portfolios:voltarget}

\begin{build}
\bfield{Purpose} Variance forecasts, target weights with caps and bands, P\&L with costs, and an aggregate-flow simulator with feedback.
\bfield{Interface} \texttt{realised\_var(r, window)}, \texttt{ewma\_var(r, com)}, \texttt{weights(var, target, cap, power)}, \texttt{banded(w, band)},
\texttt{run(r, w, cost)}, \texttt{spike\_flows(base, share, turnover, impact, target, com, cap)}.
\bfield{Rules} Forecasts from returns up to the previous close; weights capped; flows traded on the day they are decided and fed back into the forecast.
\bfield{Acceptance tests} \texttt{code/firm/voltarget/tests/}: forecasts on a constant series, weights, the band and P\&L by hand, selling after a fall
and a deeper fall with feedback.
\bfield{Stretch} Several fund types; bands; intraday timing of the flows.
\end{build}

\begin{omsources}
\item A.~Moreira and T.~Muir, ``Volatility-managed portfolios'', \emph{Journal of Finance} 72(4), 2017.
\item S.~Cederburg, M.~S.~O'Doherty, F.~Wang and X.~S.~Yan, ``On the performance of volatility-managed portfolios'', \emph{Journal of Financial
Economics} 138(1), 2020.
\item C.~R.~Harvey, E.~Hoyle, R.~Korgaonkar, S.~Rattray, M.~Sargaison and O.~van Hemert, ``The impact of volatility targeting'', \emph{Journal of
Portfolio Management} 45(1), 2018.
\end{omsources}

\section{Exercises}

\begin{exercise}[$\star$]\label{exo:s1:volatility-targeting-and-risk-managed-portfolios:1}
A fund targets 10\% volatility. What is its exposure when the forecast is 12\%? 40\%?
\end{exercise}

\begin{exercise}[$\star$]\label{exo:s1:volatility-targeting-and-risk-managed-portfolios:2}
Funds holding 0.2\% of a market's value cut their exposure from 0.73 to 0.41 in a day; the market trades 0.4\% of its value a day. What share of the
day's volume do they sell?
\end{exercise}

\begin{exercise}[$\star$]\label{exo:s1:volatility-targeting-and-risk-managed-portfolios:3}
Why does inverse-variance scaling trade more than inverse-volatility scaling?
\end{exercise}

\begin{exercise}[$\star\star$]\label{exo:s1:volatility-targeting-and-risk-managed-portfolios:4}
The synthetic market's true Sharpe ratio is $0.06/0.16$. What is it, what is the standard error of a ten-year estimate, and of the mean gain over
twenty samples when the gains have a standard deviation of 0.22?
\end{exercise}

\begin{exercise}[$\star\star$]\label{exo:s1:volatility-targeting-and-risk-managed-portfolios:5}
Why does targeting help the Sharpe ratio of equities and not of bonds in Harvey and co-authors' data?
\end{exercise}

\begin{exercise}[$\star\star$]\label{exo:s1:volatility-targeting-and-risk-managed-portfolios:6}
Explain the feedback loop between targeting funds' selling and their own volatility forecasts.
\end{exercise}

\begin{exercise}[$\star\star\star$]\label{exo:s1:volatility-targeting-and-risk-managed-portfolios:7}
\emph{Coding.} Run \texttt{spike\_flows} on \texttt{base\_path()} with a share of 0.5\% and an impact of 0.3. What happens before the spike, and why?
\end{exercise}

\begin{exercise}[$\star\star\star$]\label{exo:s1:volatility-targeting-and-risk-managed-portfolios:8}
\emph{Find the flaw.} ``Over 1990--2020 the volatility-managed market had an alpha of 4\% against the market in a spanning regression, so we will run it.''
\end{exercise}

\section{Problem: Mechanical Selling}

\begin{problem}[{Weekend problem --- targeting and its footprint}]\label{pb:s1:volatility-targeting-and-risk-managed-portfolios:1}
The chapter's synthetic markets and the public record.

\medskip\noindent\textbf{Part I --- Mechanics.}
\begin{enumerate}
\item Define a \omterm{def:s1:volatility-targeting-and-risk-managed-portfolios:managed}{volatility-managed portfolio} and a \omterm{def:s1:volatility-targeting-and-risk-managed-portfolios:flow}{leverage rebalancing flow}.
\item List the ingredients of a volatility-targeting rule.
\item What is Moreira and Muir's argument?
\item What did Cederburg and co-authors find?
\end{enumerate}

\medskip\noindent\textbf{Part II --- Sharpe ratios.}
\begin{enumerate}[resume]
\item Describe the twenty-sample experiment.
\item Give the mean Sharpe ratios and the number of samples improved.
\item Why is a real gain invisible in one decade?
\item What did Harvey and co-authors find, and what does the synthetic futures universe show?
\end{enumerate}

\medskip\noindent\textbf{Part III --- Flows.}
\begin{enumerate}[resume]
\item Describe the flow simulation and its assumptions.
\item Give the selling the day after the fall and over twenty days.
\item What does feedback do to the fall?
\item Why is the selling predictable?
\end{enumerate}

\medskip\noindent\textbf{Part IV --- The verdict.}
\begin{enumerate}[resume]
\item State the \emph{named result}: the Sharpe improvement from volatility management and the aggregate flow as a fraction of volume after a spike.
\item What does targeting reliably do, if not raise the Sharpe ratio?
\item How would you backtest a \omterm{def:s1:volatility-targeting-and-risk-managed-portfolios:managed}{volatility-managed portfolio} honestly?
\item Which strategy file profits from the others?
\item How does this chapter relate to chapter~19's portfolio targeting?
\item What would you need to size the deleveraging flow in a real market?
\item What makes the flow dangerous to the funds themselves?
\item In one sentence: what does a volatility target buy, and what does it cost the market?
\end{enumerate}
\end{problem}

\section{Interview questions}

\begin{interviewq}[$\star$ \iqroles{researcher}]\label{iq:s1:volatility-targeting-and-risk-managed-portfolios:1}
Why would scaling exposure by inverse volatility raise a Sharpe ratio?
\end{interviewq}

\begin{interviewq}[$\star\star$ \iqroles{researcher}]\label{iq:s1:volatility-targeting-and-risk-managed-portfolios:2}
What is wrong with judging a \omterm{def:s1:volatility-targeting-and-risk-managed-portfolios:managed}{volatility-managed portfolio} by its alpha in a spanning regression?
\end{interviewq}

\begin{interviewq}[$\star\star$ \iqroles{trader}]\label{iq:s1:volatility-targeting-and-risk-managed-portfolios:3}
Volatility doubled today. Which market participants must sell tomorrow, and how would you estimate how much?
\end{interviewq}

\begin{interviewq}[$\star\star$ \iqroles{risk}]\label{iq:s1:volatility-targeting-and-risk-managed-portfolios:4}
What risks does a volatility target reduce, and which does it add?
\end{interviewq}

\begin{interviewq}[$\star\star$ \iqroles{developer}]\label{iq:s1:volatility-targeting-and-risk-managed-portfolios:5}
Design the daily process that sets a volatility-targeted fund's exposure.
\end{interviewq}

\begin{interviewq}[$\star\star\star$ \iqroles{researcher}]\label{iq:s1:volatility-targeting-and-risk-managed-portfolios:6}
Returns are $r_t = \mu + \sigma_t\varepsilon_t$ with a constant $\mu$ and $\sigma_t$ known in advance. Show that the position $w_t = c/\sigma_t$ has a
higher Sharpe ratio than a constant position unless $\sigma_t$ is constant, and compute the ratio of the two in terms of the moments of $1/\sigma_t$.
\end{interviewq}
